\usepackage[spanish,es-tabla]{babel}
\usepackage[top=2cm,bottom=2cm,right=2cm,left=2cm]{geometry}
\usepackage{amsmath}
\usepackage{amssymb}
\usepackage{amsthm}
\usepackage{empheq}
\usepackage{mdframed}
\usepackage{booktabs}
\usepackage{lipsum}
\usepackage{graphicx}
\usepackage{color}
\usepackage{psfrag}
\usepackage{pgfplots}
\usepackage{bm}
\usepackage{multicol}
\usepackage{wrapfig}
\usepackage{float}
\usepackage{wrapfig}
\usepackage{multirow}
\usepackage{biblatex}
\usepackage{csquotes}
\usepackage{parskip}
\usepackage{cancel}
\usepackage{listing}
\usepackage{tikz}
\usepackage{verbatim}
\usepackage{subcaption}
\usepackage{enumitem}
\usepackage{xcolor}
\usepackage{xfrac}
\usepackage[load-configurations = abbreviations]{siunitx}
\usepackage{fix-cm}
\usepackage{yhmath}
\usepackage[bookmarks=true,colorlinks=true,linkcolor=black,citecolor=black,menucolor=black,urlcolor=black,hidelinks,breaklinks=true]{hyperref}
\usepackage{mathrsfs}
\usepackage{tcolorbox}
\usepackage{amsfonts}
\usepackage[bb=ams]{mathalpha}
\usepackage{mathrsfs}
\usepackage{xstring}

\graphicspath{{fig/}}
\addbibresource{biblio.bib}
\usepgfplotslibrary{colorbrewer}
\pgfplotsset{width=8cm,compat=1.9}
\definecolor{negro}{RGB}{0,0,0}
\definecolor{azul}{RGB}{0,0,255}
\definecolor{verde}{RGB}{0,255,0}
\definecolor{cian}{RGB}{0,255,255}
\definecolor{rojo}{RGB}{255,0,0}
\definecolor{magenta}{RGB}{255,0,255}
\definecolor{amarillo}{RGB}{255,255,0}
\definecolor{blanco}{RGB}{255,255,255}

\definecolor{limegreen}{cmyk}{0.75,0,1,0}
\definecolor{orange}{cmyk}{0,0.8,0.95,0}
\definecolor{purple}{cmyk}{0.75,1,0,0}

\newcommand{\uvec}[1]{~ \boldsymbol{\hat{\textbf{#1}}}}

\newsavebox\foobox
\newlength{\foodim}
\newcommand{\slantbox}[2][0]{\mbox{%
        \sbox{\foobox}{#2}%
        \foodim=#1\wd\foobox
        \hskip \wd\foobox
        \hskip -0.5\foodim
        \pdfsave
        \pdfsetmatrix{1 0 #1 1}%
        \llap{\usebox{\foobox}}%
        \pdfrestore
        \hskip 0.5\foodim
}}
\def\Laplacesymb{\slantbox[-.45]{$\mathscr{L}$}}
\def\inverseLaplacesymb{\Laplacesymb^{-1}}
\def\Fouriersymb{\slantbox[-.45]{$\mathscr{F}$}}
\def\Fouriersinsymb{\Fouriersymb_{\textup{s}}}
\def\Fouriercossymb{\Fouriersymb_{\textup{c}}}
\def\inverseFouriersymb{\Fouriersymb^{-1}}

\newcommand{\Laplace}[2]{
    \IfEqCase{#1}{
        {}{\Laplacesymb\left\lbrace#2\right\rbrace}
        {i}{\inverseLaplacesymb\left\lbrace#2\right\rbrace}
    }[\PackageError{tree}{Undefined option to tree: #1}{}]
}
\newcommand{\Fourier}[2]{
    \IfEqCase{#1}{
        {}{\Fouriersymb\left\lbrace#2\right\rbrace}
        {i}{\inverseFouriersymb\left\lbrace#2\right\rbrace}
        {sin}{\Fouriersinsymb\left\lbrace#2\right\rbrace}
        {cos}{\Fouriercossymb\left\lbrace#2\right\rbrace}
    }[\PackageError{tree}{Undefined option to tree: #1}{}]
}

\sisetup{
  detect-family,
  per-mode = symbol,
  separate-uncertainty = true,
  inter-unit-product = \cdot,
  %output-decimal-marker = {,},
  %exponent-product = \cdot
}

\begin{document}

\onecolumn\begin{@twocolumntrue}
    \begin{minipage}{0.3\textwidth}{\includegraphics[scale=0.24]{Escudo_UD.pdf}} % Cambiar escudo
    \end{minipage}
    \vspace{10pt}
    \begin{minipage}{0.677\textwidth}
        \begin{center}
            \vspace{12mm}
        
            \Large{\textbf{<<Título>>}}
            \vspace{3mm}
            
            \large{\textbf{<<Autores>>}}
            \vspace{2mm}
    
            \large{\textit{<<Carrera>>, <<Universidad>>}}
            \vspace{1mm}
            
            <<Fecha>> % Formato (mes) de (año)
        \end{center}
\vspace{5pt}
\end{minipage}

\centerline{\rule{0.95\textwidth}{0.4pt}} % Línea negra
\end{@twocolumntrue}

\begin{multicols}{2}

\section{Serie de Fourier}

\begin{enumerate}[leftmargin=15pt]
    \item Utilizar la calculadora \textit{DESMOS} para calcular la serie de Fourier de la función $F(t) = \ln(t)$ con $F(t+T_f) = F(t)$ si $T_f = \pi$ (Figura \ref{figenu:1}). Representar gráficamente la solución junto con su serie de Fourier para, mínimo, 5 términos.
    
    \begin{figure}[H]
    \begin{center}
    \includegraphics[scale=0.5]{fig1.png}
    \caption{$F(t) = \ln(t)$ con $F(t+T_f) = F(t)$ si $T_f = \pi$}
    \label{figenu:1}
    \end{center}
    \end{figure}
\end{enumerate}

\subsubsection*{Solución:}

Usando la calculadora \textit{DESMOS} se puede graficar la serie de Fourier para la función $F(t) = \ln(t)$. La cual es puede consultar en el siguiente link: \href{https://www.desmos.com/calculator/pu9ukvn7qk}{\textcolor{rojo}{\textit{gráfica de logaritmo natural en DESMOS}}}.

A continuación la gráfica de la serie de Fourier para $F(t)$:

\begin{figure}[H]
\begin{center}
\includegraphics[scale=0.28]{desmos-graph.png}
\caption{Gráfica de la serie de Fourier para $F(t)=\ln(t)$ evaluada hasta $n=20$}
\label{fig1}
\end{center}
\end{figure}

\begin{enumerate}[leftmargin=15pt, resume]
    \item Demostrar que la serie de Fourier para la función $F(t) = t^{2}$, en el intervalo $(-\pi,~\pi)$, y con $F(t) = F(t + T_f)$ es:
    \begin{equation*}
        F(t) = \frac{1}{3} \pi^{2} + 4 \sum_{n = 1}^{\infty} \frac{(-1)^{n}}{n^{2}} \cos(nt)
    \end{equation*}
    Luego, en la serie de Fourier, haga $t = \pi$, y demostrar que:
    \begin{equation*}
        \sum_{1}^{\infty} \frac{1}{n^{2}} = 1 + \frac{1}{4} + \frac{1}{9} + \frac{1}{16} + \frac{1}{25} + \cdots = \frac{\pi^{2}}{6}
    \end{equation*}
\end{enumerate}

\subsubsection*{Solución:}

Primeramente, es sencillo analizar la función dada si fuese una función par o impar ya que para $F(t) = t^{2}$, $F(t)$ siempre es una función par, por lo tanto, los coeficientes que se van a calcular son:
\begin{align}
    a_0 &= \frac{4}{T_f} \int\limits_{0}^{\sfrac{T_{f}}{2}} F(t) \,dt \label{pro2:eq1}\\
    a_n &= \frac{4}{T_f} \int\limits_{0}^{\sfrac{T_{f}}{2}} F(t) \cos(n \omega_f t) \,dt\label{pro2:eq2}
\end{align}
Y $b_n = 0$, para $\omega_f = \sfrac{2 \pi}{T_f}$.

Conociendo esto, podemos comenzar a calcular los coeficientes para la serie de Fourier, por tanto:

Para $a_0$ con $T_f = 2 \pi$ y $\omega_f = 1$:
\begin{align*}
    a_0 &= \frac{\cancel{4}}{\cancel{2} \pi} \int\limits_{0}^{\sfrac{\cancel{2} \pi}{\cancel{2}}} t^{2} \,dt\\
    &= \frac{2}{\pi} \int\limits_{0}^{\pi} t^{2} \,dt\\
    &= \frac{2}{\pi} \left(\left.\frac{t^{3}}{3} \right|_{0}^{\pi}\right)\\
    &= \frac{2}{\pi} \left(\frac{\left(\pi\right)^{3}}{3} - \cancel{\frac{\left(0\right)^{3}}{3}}\right)\\
    &= \frac{2}{\cancel{\pi}} \left(\frac{\pi^{\cancel{3}}}{3}\right)\\
    &= \frac{2\pi^{2}}{3}
\end{align*}
Por tanto, el coeficiente $a_0$ es:
\begin{equation}
    a_0 = \frac{2\pi^{2}}{3} \label{pro2:res1}
\end{equation}
Para $a_n$ con $T_f = 2 \pi$ y $\omega_f = 1$:
\begin{align*}
    a_n &= \frac{\cancel{4}}{\cancel{2} \pi} \int\limits_{0}^{\sfrac{\cancel{2} \pi}{\cancel{2}}} t^{2} \cos(nt) \,dt\\
    &= \frac{2}{\pi} \int\limits_{0}^{\pi} t^{2} \cos(nt) \,dt\\
\end{align*}
Sustituyendo al usar integración por partes:
\begin{align*}
    u&=t^{2}           &  dv&=\cos\left(nt\right) dt\\
    du&=2tdt         &  v&=\frac{1}{n} \sin\left(nt\right)
\end{align*}
Entonces:
\begin{align*}
    a_n = \frac{2}{\pi} \left[\frac{t^{2}}{n}\sin(nt) - \frac{2}{n} \int t \sin(nt) \,dt \right]_{0}^{\pi}
\end{align*}
Sustituyendo al usar integración por partes, nuevamente:
\begin{align*}
    u&=t          &  dv&=\sin\left(nt\right) dt\\
    du&=dt         &  v&=-\frac{1}{n} \cos\left(nt\right)
\end{align*}
Entonces:
\begin{align*}
    a_n =& \frac{2}{\pi} \left[\frac{t^{2}}{n}\sin(nt)\right]_{0}^{\pi} - \frac{2}{n}\\
    &\left[-\frac{t}{n}\cos(nt) + \frac{1}{n} \int \cos(nt) \,dt \right]_{0}^{\pi}\\
    =& \frac{2}{\pi} \left[\frac{t^{2}}{n}\sin(nt)\right]_{0}^{\pi} - \frac{2}{n}\\ 
    &\left[-\frac{t}{n}\cos(nt) + \frac{1}{n} \left(\frac{1}{n} \sin(nt)\right) \right]_{0}^{\pi}\\
    =& \frac{2}{\pi} \left[\frac{t^{2}}{n} \sin(nt) + \frac{2t}{n^{2}} \cos(nt) - \frac{2}{n^{3}} \sin(nt)\right]_{0}^{\pi}\\
    =& \frac{2}{\pi} \left(\frac{\pi^{2}}{n} \sin(n\pi) + \frac{2\pi}{n^{2}} \cos(n\pi) - \frac{2}{n^{3}} \sin(n\pi)\right)\\
    &\textup{Evaluando en 0, el argumento se cancela}
\end{align*}
Ya que $\sin(n\pi) = 0$ y $\cos(n\pi) = (-1)^{n}$, resulta finalmente:
\begin{align*}
    a_n &= \frac{2}{\pi} \left(\cancel{\frac{\pi^{2}}{n}(0)} + \frac{2\pi}{n^{2}}\left(-1\right)^{n} - \cancel{\frac{2}{n^{3}}(0)}\right)\\
    &= \frac{2}{\cancel{\pi}} \left(\frac{2\cancel{\pi}}{n^{2}}\left(-1\right)^{n}\right)\\
    &= 4 \frac{\left(-1\right)^{n}}{n^{2}}
\end{align*}
Por tanto, el coeficiente $a_n$ es:
\begin{equation}
    a_n = 4 \frac{\left(-1\right)^{n}}{n^{2}} \label{pro2:res2}
\end{equation}
Conociendo que $b_n = 0$, usando (\ref{pro2:res1}) y (\ref{pro2:res2}), la serie trigonométrica de Fourier para $F(t) = F(t + 2\pi) = t^{2}$ es:
\begin{equation*}
    F(t) = \frac{1}{\cancel{2}}\left(\frac{\cancel{2}\pi^{2}}{3}\right) + \sum_{n = 1}^{\infty} 4 \frac{\left(-1\right)^{n}}{n^{2}} \cos(nt)
\end{equation*}
Desarrollando, resulta:
\begin{equation}
    \textcolor{rojo}{F(t) = \frac{1}{3}\pi^{2} + 4 \sum_{n = 1}^{\infty} \frac{\left(-1\right)^{n}}{n^{2}} \cos(nt)} \label{pro2:res3}
\end{equation}
Como el ejercicio pide evaluar la serie de Fourier obtenida, evaluando $t = \pi$ en (\ref{pro2:res3}):
\begin{align*}
    (\pi)^{2} &= \frac{1}{3}\pi^{2} + 4 \sum_{n = 1}^{\infty} \frac{\left(-1\right)^{n}}{n^{2}} \cos(n\pi)\\
    &\textup{Ya que } \cos(n\pi) = (-1)^{n}\\
    &\pi^{2} = \frac{1}{3}\pi^{2} + 4 \sum_{n = 1}^{\infty} \frac{\left(-1\right)^{2n}}{n^{2}}
\end{align*}
Como la forma par de $n$ es $2n$, entonces si $(-1)$ se eleva a un valor número par el resultado es siempre 1, es decir, $(-1)^{2n} = 1$, reemplazando esto y continuando con el desarrollo:
\begin{align*}
    \pi^{2} &= \frac{1}{3}\pi^{2} + 4 \sum_{n = 1}^{\infty} \frac{1}{n^{2}}\\
    \pi^{2} - \frac{1}{3}\pi^{2} &= 4 \sum_{n = 1}^{\infty} \frac{1}{n^{2}}\\
    \pi^{2}\left(1 - \frac{1}{3}\right) &= 4 \sum_{n = 1}^{\infty} \frac{1}{n^{2}}\\
    \pi^{2}\left(\frac{2}{3}\right) &= 4 \sum_{n = 1}^{\infty} \frac{1}{n^{2}}\\
    \sum_{n = 1}^{\infty} \frac{1}{n^{2}} &= \frac{\pi^{2}}{\cancel{4}}\left(\frac{\cancel{2}}{3}\right)\\
    \sum_{n = 1}^{\infty} \frac{1}{n^{2}} &= \frac{\pi^{2}}{6}
\end{align*}
De esta manera, si evaluamos en la sumatoria para los primeros 5 términos de $n$, el resultado final es:
\begin{equation}
    \textcolor{rojo}{\sum_{n = 1}^{\infty} \frac{1}{n^{2}} = 1 + \frac{1}{4} + \frac{1}{9} + \frac{1}{16} + \frac{1}{25} + \cdots = \frac{\pi^{2}}{6}} \label{pro2:res4}
\end{equation}

\begin{enumerate}[leftmargin=15pt, resume]
    \item Integrar la serie de Fourier del problema anterior: $F(t) = t^{2}$, para obtener:
    \begin{equation*}
        \sum_{1}^{\infty} (-1)^{n} \frac{\sin(nt)}{n^{3}} = \frac{1}{12} t\left(t^{2} - \pi^{2}\right)
    \end{equation*}
    y
    \begin{equation*}
        \sum_{1}^{\infty} \frac{1}{n^{6}} = \frac{\pi^{6}}{945}
    \end{equation*}
\end{enumerate}

\subsubsection*{Solución:}

Tomando la serie de Fourier del problema anterior obtenida en (\ref{pro2:res3}), si la integramos indefinidamente en ambas partes con respecto a $t$, haciendo el respectivo tratamiento:
\begin{equation*}
    \int t^{2} \,dt = \int \left(\frac{1}{3}\pi^{2} + 4 \sum_{n = 1}^{\infty} \frac{\left(-1\right)^{n}}{n^{2}} \cos(nt)\right) \,dt
\end{equation*}
Ya que la integral de una suma es la suma de sus integrales, esto gracias a la linealidad de la integral, se puede sacar la sumatoria de la integral y solamente integrar la parte que dependa netamente de $t$, mismo caso para las constantes distribuidas. 

Resolviendo la integral resulta:
\begin{align*}
    \int t^{2} \,dt &= \frac{1}{3}\pi^{2 }\int \,dt + 4 \sum_{n = 1}^{\infty} \frac{\left(-1\right)^{n}}{n^{2}} \int \cos(nt) \,dt\\
    \frac{t^{3}}{3} &= \frac{1}{3}\pi^{2} \left(t\right) + 4 \sum_{n = 1}^{\infty} \frac{\left(-1\right)^{n}}{n^{2}} \left(\frac{1}{n}\sin(nt)\right)\\
    \frac{t^{3}}{3} &= \frac{t}{3}\pi^{2} + 4 \sum_{n = 1}^{\infty} \frac{\left(-1\right)^{n}}{n^{3}} \sin(nt)\\
    \frac{t^{3}}{3} - \frac{t}{3}\pi^{2} &= 4 \sum_{n = 1}^{\infty} \frac{\left(-1\right)^{n}}{n^{3}} \sin(nt)\\
    \frac{t}{12} \left(t^{2} - \pi^{2}\right) &= \sum_{n = 1}^{\infty} \frac{\left(-1\right)^{n}}{n^{3}} \sin(nt)
\end{align*}
Por tanto, la integral de la serie de Fourier de $F(t) = t^{2}$ es:
\begin{equation}
    \textcolor{rojo}{\sum_{n = 1}^{\infty} \frac{\left(-1\right)^{n}}{n^{3}} \sin(nt) = \frac{t}{12} \left(t^{2} - \pi^{2}\right)} \label{pro3:res1}
\end{equation}
Para la segunda parte de la demostración que exige el punto, es necesario aplicar la \textbf{identidad de Parseval} la cual viene dada por:
\begin{equation*}
    \frac{2}{T_f} \int\limits_{-\sfrac{T_f}{2}}^{\sfrac{T_f}{2}} \left[F(t)\right]^{2} \,dt = \frac{\left(a_0\right)^{2}}{2} + \sum_{n=1}^{\infty} \left[\left(a_n\right)^{2} + \left(b_n\right)^{2}\right]
\end{equation*}
La identidad mostrada se va a aplicar sobre la parte dependiente de $t$ obtenida en (\ref{pro3:res1}), es decir, $F(t) = \frac{t}{12} \left(t^{2} - \pi^{2}\right)$, que si analizamos el resultado de la integral, tiene la forma de una serie de Fourier, efectivamente se podría suponer que la integral de una serie de Fourier es igual a otra serie de Fourier, por lo que si recordamos la forma de una serie de Fourier, sabemos que el coeficiente que acompaña a la función $\sin(n\omega_ft)$ es $b_n$, por lo que se puede deducir también que $a_0 = a_n = 0$ y $b_n = \frac{(-1)^{n}}{n^{3}}$, donde gracias al argumento del seno ($\sin(nt)$), se ve claramente que $\omega_f = 1$ lo que implica que $T_f = 2\pi$.

Usando todos estos resultados deducidos de solo analizar lo obtenido (\ref{pro3:res1}) en la identidad, resulta:
\begin{align*}
    \frac{\cancel{2}}{\cancel{2}\pi} \int\limits_{-\sfrac{\cancel{2}\pi}{\cancel{2}}}^{\sfrac{\cancel{2}\pi}{\cancel{2}}} \left[\frac{t}{12} \left(t^{2} - \pi^{2}\right)\right]^{2} \,dt &= \cancel{\frac{(0)^{2}}{2}} + \sum_{n=1}^{\infty} \left(\frac{(-1)^{n}}{n^{3}}\right)^{2}\\
    \frac{1}{\pi} \int\limits_{-\pi}^{\pi} \left(\frac{t^{3} - t\pi^{2}}{12}\right)^{2} \,dt &= \sum_{n=1}^{\infty} \frac{(-1)^{2n}}{n^{6}}
\end{align*}
Se va a desarrollar a continuación solamente la parte de la izquierda correspondiente a la integral, esto debido a que la parte de la derecha se va a mantener constante y la integral será algo compleja de desarrollar. Siguiendo el mismo planteamiento del punto anterior dónde $(-1)^{2n} = 1$ y realizando lo mencionado:

La integral resulta:
\begin{gather*}
    \int\limits_{-\pi}^{\pi} \left(\frac{t^{3} - t\pi^{2}}{12}\right)^{2} \,dt\\
    \int\limits_{-\pi}^{\pi} \frac{t^{6} - 2t^{4}\pi^{2} + t^{2}\pi^{4}}{144} \,dt\\
    \frac{1}{144} \left[\int\limits_{-\pi}^{\pi} t^{6} \,dt - 2\pi^{2}\int\limits_{-\pi}^{\pi} t^{4} \,dt + \pi^{4}\int\limits_{-\pi}^{\pi} t^{2} \,dt\right]\\
    \frac{1}{144} \left[\left(\left.\frac{t^{7}}{7}\right|_{-\pi}^{\pi}\right) - 2\pi^{2}\left(\left.\frac{t^{5}}{5}\right|_{-\pi}^{\pi}\right) + \pi^{4}\left(\left.\frac{t^{3}}{3}\right|_{-\pi}^{\pi}\right)\right]
\end{gather*}
Desarrollando las respectivas evaluaciones y multiplicando por los factores de cada integral:
\begin{gather*}
    \frac{1}{144} \left[\frac{2\pi^{7}}{7} - \frac{4\pi^{7}}{5} + \frac{2\pi^{7}}{3}\right]\\
    \frac{1}{144} \left[\frac{30\pi^{7} - 84\pi^{7} + 70\pi^{7}}{105}\right]\\
    \frac{1}{\cancel{144}} \left[\frac{\cancel{16}\pi^{7}}{105}\right]\\
    \frac{1}{9} \left[\frac{\pi^{7}}{105}\right]\\
    \frac{\pi^{7}}{945}
\end{gather*}
Por tanto el valor de la integral es:
\begin{equation}
    \int\limits_{-\pi}^{\pi} \left(\frac{t^{3} - t\pi^{2}}{12}\right)^{2} \,dt = \frac{\pi^{7}}{945} \label{pro3:res2}
\end{equation}
Reemplazando el valor obtenido en (\ref{pro3:res2}) en el tratamiento que ya se estaba realizando con base en la identidad de Parseval:
\begin{align*}
    \frac{1}{\cancel{\pi}} \left(\frac{\pi^{\cancel{7}}}{945}\right) &= \sum_{n=1}^{\infty} \frac{1}{n^{6}}\\
    \frac{\pi^{6}}{945} &= \sum_{n=1}^{\infty} \frac{1}{n^{6}}
\end{align*}
Por tanto, la identidad de Parseval aplicada para (\ref{pro3:res1}) es:
\begin{equation}
    \textcolor{rojo}{\sum_{n=1}^{\infty} \frac{1}{n^{6}} = \frac{\pi^{6}}{945}} \label{pro3:res3}
\end{equation}

\begin{enumerate}[leftmargin=15pt, resume]
    \item Utilizar el teorema de Parseval para probar que:
    \begin{equation*}
        \sum_{1}^{\infty} \frac{1}{(2n-1)} = \frac{\pi^{2}}{8}
    \end{equation*}
\end{enumerate}

\subsubsection*{Solución:}

Para la solución de este problema se va a volver a aplicar la identidad y/o teorema de Parseval, que recordando su forma:
\begin{equation*}
    \frac{2}{T_f} \int\limits_{-\sfrac{T_f}{2}}^{\sfrac{T_f}{2}} \left[F(t)\right]^{2} \,dt = \frac{\left(a_0\right)^{2}}{2} + \sum_{n=1}^{\infty} \left[\left(a_n\right)^{2} + \left(b_n\right)^{2}\right]
\end{equation*}
Se va a aplicar el teorema para la función periódica a trozos:
\begin{equation}
    F(t) =
    \begin{cases}
    -1,&\text{si } \sfrac{-\pi}{2} \leq t < 0\\
    1,&\text{si } \hspace{18pt} 0 \leq t < \sfrac{\pi}{2}
    \end{cases}
    \label{pro4:eq1}
\end{equation}
Que si recordamos de la anterior tarea (\textit{``Tarea 1: funciones periódicas y series trigonométricas de Fourier''}) realizada por el mismo autor de esta tarea, ya se sabe que la serie de Fourier de esta función viene expresada de la forma:
\begin{equation}
    F(t) = \sum_{\mathrm{impar}}^{\infty} \frac{4}{n\pi} \sin(2nt) \label{pro4:eq2}
\end{equation}
Donde $a_0 = a_n = 0$, $b_n = \frac{4}{n\pi}$ y $T_f = \pi$. Hay que tener en cuenta que la $n$ utilizada representa un valor de $n$ \textbf{impar}, esto es importante ya que si no se especifica ese pequeño detalle el resultado puede llegar a ser completamente diferente.

Desarrollando el teorema de Parseval para (\ref{pro4:eq1}), resulta:
\begin{align*}
    \frac{2}{\pi} \int\limits_{-\sfrac{\pi}{2}}^{\sfrac{\pi}{2}} \left[F(t)\right]^{2} \,dt &= \cancel{\frac{\left(0\right)^{2}}{2}} + \sum_{n=1}^{\infty} \left(\frac{4}{n\pi}\right)^{2}\\
    \frac{2}{\pi} \left[\int\limits_{-\sfrac{\pi}{2}}^{0} \left[-1\right]^{2} \,dt + \int\limits_{0}^{\sfrac{\pi}{2}} \left[1\right]^{2} \,dt\right] &= \sum_{n=1}^{\infty} \frac{16}{n^{2}\pi^{2}}\\
    \frac{\cancel{2}}{\cancel{\pi}} \left[\int\limits_{-\sfrac{\pi}{2}}^{0} \,dt + \int\limits_{0}^{\sfrac{\pi}{2}} \,dt\right] &= \frac{\cancel{16}}{\pi^{\cancel{2}}} \sum_{n=1}^{\infty} \frac{1}{n^{2}}\\
    \left(\left.t\right|_{\sfrac{-\pi}{2}}^{0}\right) + \left(\left.t \right|_{0}^{\sfrac{\pi}{2}}\right) &= \frac{8}{\pi} \sum_{n=1}^{\infty} \frac{1}{n^{2}}\\
    \left(\frac{\pi}{2}\right) + \left(\frac{\pi}{2}\right) &= \frac{8}{\pi} \sum_{n=1}^{\infty} \frac{1}{n^{2}}\\
    \pi &= \frac{8}{\pi} \sum_{n=1}^{\infty} \frac{1}{n^{2}}\\
    \frac{\pi^{2}}{8} &= \sum_{n=1}^{\infty} \frac{1}{n^{2}}
\end{align*}
Como se había mencionado previamente, el valor de $n$ representa a una $n$ la cual es impar, o mejor dicho, $n$ es impar, lo que implica que puede ser reescrito de la forma $n = 2n-1$. 

Sin embargo, para nuestro caso particular dónde tenemos $n^{2}$ en vez de simplemente $n$, esto puede ser nuevamente reescrito sencillamente como un número impar general, esto gracias a que si $n$ es impar, $n^{2}$ también lo es; a continuación dejo un video con la respectiva demostración de este hecho: \href{https://youtu.be/XMB3Aa6SscQ}{\textit{si $n$ es impar entonces $n^{2}$ es impar}}.

Por tanto, el teorema de Parseval aplicado a (\ref{pro4:eq1}) es:
\begin{equation}
    \textcolor{rojo}{\sum_{n=1}^{\infty} \frac{1}{2n-1} = \frac{\pi^{2}}{8}} \label{pro4:res1}
\end{equation}

\section{Ecuaciones diferenciales}

\begin{enumerate}[leftmargin=15pt]
    \item Resolver la siguiente ecuación diferencial:
    \begin{equation*}
        \frac{d^{2}x}{dt^{2}} + \lambda \frac{dx}{dt} + \omega_{0}^{2}x = \frac{1}{m} F(t)
    \end{equation*}
    Para los siguientes casos:
    \begin{enumerate}
        \item Caso 1:
        \begin{equation*}
            F(t) =
            \begin{cases}
            -1,&\text{si } \sfrac{-\pi}{2} \leq t < 0\\
            1,&\text{si } \hspace{18pt} 0 \leq t < \sfrac{\pi}{2}
            \end{cases}
        \end{equation*}
        \begin{center}$F(t)=F(t+\pi)$\end{center}
        \item Caso 2:
        \begin{equation*}
            F(t) =
            \begin{cases}
            1 + \frac{4t}{\pi},&\text{si } \sfrac{-\pi}{2} \leq t < 0\\
            1 - \frac{4t}{\pi},&\text{si } \hspace{18pt} 0 \leq t < \sfrac{\pi}{2}
            \end{cases}
        \end{equation*}
        \begin{center}$F(t)=F(t+\pi)$\end{center}
    \end{enumerate}
\end{enumerate}

\subsubsection*{Solución:}

Sabiendo que la solución de cualquier ecuación diferencial lineal se compone de la superposición de dos soluciones, una complementaria y una particular. La solución complementaria se halla tras resolver la ecuación diferencial homogénea asociada a la ecuación diferencial general, para el caso de una ecuación diferencial que de por sí es homogénea, pues la solución general termina siendo solamente la complementaria, cosa que para nuestro caso no es así. La solución particular se hallará por el método de los coeficientes indeterminados, esto depende del valor de $F(t)$ que más adelante se analizará para los casos correspondientes.

A continuación, se hallará la solución complementaria tomando la ecuación diferencial homogénea asociada a la ecuación diferencial general, la cuál es:
\begin{equation}
    \frac{d^{2}x}{dt^{2}} + \lambda \frac{dx}{dt} + \omega_{0}^{2}x = 0 \label{pro5:eq1}
\end{equation}
Donde su desarrollo va a ser a partir de la parte compleja de la ecuación diferencial.

Si nosotros sabemos que una función compleja se compone de dos partes, una parte real equivalente a una función real, y una parte imaginaria equivalente a una función real multiplicada por $i$,
\begin{equation*}
    Z(t) = X(t) + i Y(t)
\end{equation*}
podemos asumir el mismo tratamiento para una ecuación diferencial, por tanto, la parte compleja de la ecuación diferencial se puede expresar como la combinación lineal de dos ecuaciones diferenciales equivalentes, usando la notación de Newton:
\begin{equation*}
\ddot{Z} + \lambda \dot{Z} + \omega_{0}^{2}Z = 0
    \begin{cases}
        &\ddot{X} + \lambda \dot{X} + \omega_{0}^{2}X = 0\\
        &\ddot{Y} + \lambda \dot{Y} + \omega_{0}^{2}Y = 0
    \end{cases}
\end{equation*}
Donde $X(t)$ es la parte real de la función compleja $Z(t)$, y $Y(t)$ corresponde a la parte imaginaria. De esta manera, es solo cuestión de reemplazar la función original (\ref{pro5:eq1}) por la nueva función compleja en notación de Newton.

Para resolver la ecuación diferencial homogénea lineal con coeficientes constantes se va a proponer un \textit{anzats} (intento en alemán), o mejor dicho, suponer una solución basándose en la forma de la ecuación diferencial, ya que estamos viendo una ecuación diferencial con funciones complejas, se supone que la solución sea la forma elemental de un número complejo, la cual es:
\begin{equation*}
    z = \rho e^{i\theta}
\end{equation*}
Donde $\rho$ y $\theta$ pertenecen a $\mathbb{R}$. De esta manera, se supone de igual forma la solución teniendo la misma forma pero ajustada a nuestro caso, por tanto:
\begin{equation}
    z = Ae^{iBt} \label{pro5:eq2}
\end{equation}
Tomando (\ref{pro5:eq2}) como anzats, se deriva dos veces buscando que sea solución de la ecuación diferencial.
\begin{align*}
    z &= Ae^{iBt}\\
    \dot{z} &= iABe^{iBt}\\
    \ddot{z} &= -A^{2}B^{2}e^{iBt}
\end{align*}
Reemplazando en la ecuación diferencial:
\begin{align*}
    \left(-A^{2}B^{2}e^{iBt}\right) + \lambda\left(iABe^{iBt}\right) + \omega_{0}^{2}\left(Ae^{iBt}\right) &= 0\\
    Ae^{iBt}\left(-B^{2} + \lambda iB + \omega_{0}^{2}\right) &= 0
\end{align*}
De donde se tienen dos posibles soluciones, una trivial la cual sería que $A = 0$, cosa que no nos daría ningún tipo de información, y la otra sería que el argumento factorizado sea igual a 0, caso en el cuál nos vamos a centrar.

Nos podemos dar cuenta que dicho argumento tiene la forma de una ecuación general cuadrática, por lo que aplicando la \textit{fórmula de Bhaskara} resulta:
\begin{align*}
    B &= \frac{-\lambda i \pm \sqrt{\left(\lambda i\right)^{2} -4\left(-1\right)\left(\omega_{0}^{2}\right)}}{-2}\\
    B &= \frac{i\lambda}{2} \mp \frac{\sqrt{-\lambda^{2} + 4^{\omega_{0}^{2}}}}{2}\\
    B &= \frac{i\lambda}{2} \mp \textcolor{azul}{\sqrt{\omega_{0}^{2} - \left(\frac{\lambda}{2}\right)^{2}}}
\end{align*}
La raíz de color azul se va a reemplazar por $\omega_a$, esto más que todo por efectos prácticos que curiosamente contiene un sentido físico. De esta manera obtenemos dos soluciones las cuales son:
\begin{equation}
    B_{1} = \frac{i \lambda}{2} - \omega_a ~~~;~~~ B_{2} = \frac{i \lambda}{2} + \omega_a \label{pro5:res1}
\end{equation}
Reemplazando las soluciones obtenidas en (\ref{pro5:res1}) en (\ref{pro5:eq2}):
\begin{align*}
    z_1 &= Ae^{i\left(\frac{i \lambda}{2} - \omega_a\right)t} & z_2 &= Ae^{i\left(\frac{i \lambda}{2} + \omega_a\right)t}\\
    z_1 &= Ae^{iB_{1}t} & z_2 &= Ae^{iB_{2}t}
\end{align*}
Por tanto, la solución general para la ecuación diferencial corresponderán a una combinación lineal de las soluciones obtenidas:
\begin{equation*}
    z(t) = \sum_{n=1}^{2} C_{n} Ae^{iB_{n}t} \Rightarrow z(t) = C_{1} Ae^{iB_{1}t} + C_{2} Ae^{iB_{2}t}
\end{equation*}
Que si tomamos $c_{1} = C_{1}A$ y $c_{2} = C_{2}A$, se puede escribir finalmente como:
\begin{equation}
    z(t) = c_{1} e^{iB_{1}t} + c_{2} e^{iB_{2}t} \label{pro5:res2}
\end{equation}
Ya teniendo la solución de la parte compleja de la ecuación diferencial, se debe hallar la parte real ya que es esa la solución que necesitamos, para ello se usará la definición de la conjugada de un número complejo dónde decimos que si sumamos un número complejo $z$ con su conjugada $\bar{z}$, el resultado es la parte real de $z$, que para nuestro caso sería $\mathrm{Re}[z(t)] = x(t)$.

Sin embargo, no solo $z$ es complejo, si no que también sus constantes arbitrarias $c_{1}$ y $c_{2}$ por lo que necesitamos igualar $\bar{z}$ con $z$ para hallar las condiciones necesarias para las constantes, entonces:
\begin{align*}
    z &= \bar{z}\\
    \cancel{e^{-\frac{\lambda}{2}}} \left(c_{1} e^{-i\omega_a t} + c_{2} e^{i\omega_a t}\right) &= \cancel{e^{-\frac{\lambda}{2}}} \left(\bar{c_{1}} e^{i\omega_a t} + \bar{c_{2}} e^{-i\omega_a t}\right)\\
    c_{1} e^{-i\omega_a t} + c_{2} e^{i\omega_a t} &= \bar{c_{1}} e^{i\omega_a t} + \bar{c_{2}} e^{-i\omega_a t}
\end{align*}
De donde se concluye que $c_{1} = \bar{c_{2}}$ y $c_{2} = \bar{c_{1}}$. Realizando un cambio de variables dándole un valor a las constantes respetando la forma general de un número complejo donde:
\begin{align*}
    c_{1} = D e^{i\alpha} &\Rightarrow \bar{c_{2}} = D e^{i\alpha}\\
    c_{2} = D e^{-i\alpha} &\Rightarrow \bar{c_{1}} = D e^{-i\alpha}
\end{align*}
Ya con esto, sumando y reemplazando los valores correspondientes:
\begin{align*}
    \mathrm{Re} = x(t) &= 2De^{-\frac{\lambda}{2}} \left(e^{i\alpha}e^{i\omega_at} + e^{-i\alpha}e^{-i\omega_at}\right)\\
    x(t) &= 2De^{-\frac{\lambda}{2}} \left(e^{i\left(\omega_at + \alpha\right)} + e^{-i\left(\omega_at + \alpha\right)}\right)\\
\end{align*}
Si llamamos a $2D = A_0$, desarrollando las exponenciales complejas, resulta finalmente:
\begin{equation}
    \textcolor{rojo}{x(t) = A_0e^{-\frac{\lambda}{2}} \cos\left(\omega_a t + \alpha_a\right)} \label{pro5:res3}
\end{equation}
Teniendo la solución complementaria de la ecuación diferencial expuesta, ahora se pueden hallar las soluciones para cada caso.

$a)$ Para el caso 1, con:
\begin{equation*}
    F(t) =
    \begin{cases}
    -1,&\text{si } \sfrac{-\pi}{2} \leq t < 0\\
    1,&\text{si } \hspace{18pt} 0 \leq t < \sfrac{\pi}{2}
    \end{cases}
\end{equation*}
\begin{center}$F(t)=F(t+\pi)$\end{center}
Que si nos remitimos nuevamente a la tarea anterior (\textit{``Tarea 1: funciones periódicas y series trigonométricas de Fourier''}), sabemos ya la serie trigonométrica de Fourier para $F(t)$ es:
\begin{equation}
    F(t) = \sum_{\mathrm{impar}}^{\infty} \frac{4}{n\pi} \sin\left(2nt\right)\label{pro5:eq3}
\end{equation}
Donde $a_0 = a_n = 0$ y $b_n = \frac{4}{n\pi}$. Por efecto prácticos es necesario expresar esta serie de la forma:
\begin{equation*}
    F(t) = c_0 + \sum_{n=1}^{\infty} c_n \cos\left(n\omega_f t - \theta_n\right) 
\end{equation*}
Donde $c_0 = \frac{a_0}{2}$, $c_n = \sqrt{a_{n}^{2} + b_{n}^{2}}$ y $\theta_n = \arctan\left(\frac{b_n}{a_n}\right)$. Que para este caso, debido a las condiciones que nos ofrece (\ref{pro5:eq3}), solo es cuestión de aplicar un desfase a la función seno de tal forma que esta sea un coseno, dicho desfase será $\theta_n = \frac{\pi}{2}$. La serie de Fourier resulta:
\begin{equation}
    F(t) = \sum_{\mathrm{impar}}^{\infty} \frac{4}{n\pi} \cos\left(2nt - \frac{\pi}{2}\right) \label{pro5:eq4}
\end{equation}
Realizando el mismo tratamiento hecho para la solución complementaria, se va a tomar la parte compleja para (\ref{pro5:eq4}) que gracias a que está expresada con un coseno, se puede rápidamente cómo:
\begin{equation}
    Z(t) = \sum_{\mathrm{impar}}^{\infty} \frac{4}{n\pi} e^{i\left(2nt - \frac{\pi}{2}\right)} \label{pro5:eq5}
\end{equation}
Una pequeña aclaración, cuando se está usando $Z(t)$ como la parte compleja o la función compleja equivalente para $F(t)$, en ningún momento se está hablando de tomar eso como si fuese una solución ya para la ecuación, es solo notación, que quizá más adelante tenga alguna equivalencia pero nuevamente, es solo notación.

Partiendo de ahí, usando el método de los coeficientes indeterminados, si tenemos una forma para $Z(t)$ así, se espera que las soluciones sean iguales o por lo menos equivalentes, por lo que se va a suponer un $Z_p$ el cual sea solución. Para demostrar que es solución, se derivará dos veces y reemplazará en la ecuación diferencial no homogénea asociada y se buscarán las condiciones para que se cumpla que sea solución. Entonces:
\begin{align*}
    Z_p &= A + \sum_{\mathrm{impar}}^{\infty} B e^{i\left(2nt - \frac{\pi}{2} - \varphi\right)}\\
    \dot{Z_p} &= \sum_{\mathrm{impar}}^{\infty} i2nB e^{i\left(2nt - \frac{\pi}{2} - \varphi\right)}\\
    \ddot{Z_p} &= \sum_{\mathrm{impar}}^{\infty} -4n^{2}B^{2} e^{i\left(2nt - \frac{\pi}{2} - \varphi\right)}
\end{align*}
Reemplazando en la ecuación diferencial, desarrollando y usando (\ref{pro5:eq5}):
\end{multicols}

\centerline{\rule{0.95\textwidth}{0.4pt}} % Línea negra
\begin{gather*}
    \sum_{\mathrm{impar}}^{\infty} -4n^{2}B^{2} e^{i\left(2nt - \frac{\pi}{2} - \varphi\right)} + \lambda\left(\sum_{\mathrm{impar}}^{\infty} i2nB e^{i\left(2nt - \frac{\pi}{2} - \varphi\right)}\right) + \omega_{0}^{2}\left(A + \sum_{\mathrm{impar}}^{\infty} B e^{i\left(2nt - \frac{\pi}{2} - \varphi\right)}\right) = \frac{1}{m} Z(t)
\end{gather*}
\begin{gather*}
    \left[\sum_{\mathrm{impar}}^{\infty} e^{i\left(2nt - \frac{\pi}{2}\right)}e^{-\varphi} \left(-4n^{2}B^{2} + \lambda i2nB + \omega_{0}^{2} B\right) \right] + \omega_{0}^{2} A = \frac{1}{m} \left(\sum_{\mathrm{impar}}^{\infty} \frac{4}{n\pi} e^{i\left(2nt - \frac{\pi}{2}\right)}\right)
\end{gather*}
\begin{center}
    Si igualamos las componentes similares, obtenemos qué:
\end{center}
\begin{gather*}
    \omega_{0}^{2} A = 0 \Longrightarrow A = 0\\
    e^{-\varphi} \left(-4n^{2}B^{2} + \lambda i2nB + \omega_{0}^{2} B\right) = \frac{4}{mn\pi}
\end{gather*}
\begin{center}
    Desarrollando lo obtenido para $B$ y para $\varphi$, resulta:
\end{center}
\begin{gather*}
    \frac{4}{mn\pi} = e^{-\varphi} \left(-4n^{2}B^{2} + \lambda i2nB + \omega_{0}^{2} B\right)\\
    \frac{4}{mn\pi} e^{\varphi} = B\left(-4n^{2} + \omega_{0}^{2}\right) + i\left(2nB\lambda\right)\\
    \frac{4}{mn\pi} \left(\cos(\varphi) + i\sin(\varphi)\right) = B\left(-4n^{2} + \omega_{0}^{2}\right) + i\left(2nB\lambda\right)
\end{gather*}
\begin{center}
    Si se iguala la parte real con la compleja en ambas partes, se genera el siguiente sistema de ecuaciones:
\end{center}
\begin{gather*}
    \begin{cases}
    \frac{4}{mn\pi} \cos(\varphi) = B\left(-4n^{2} + \omega_{0}^{2}\right) & (1)\\
    \frac{4}{mn\pi} \sin(\varphi) = 2nB\lambda & (2)
    \end{cases}
\end{gather*}
\begin{center}
    Usando (2) en (1):
\end{center}
\begin{gather*}
    \frac{\cancel{\frac{4}{mn\pi}}}{\cancel{\frac{4}{mn\pi}}} \tan(\varphi) = \frac{2n\cancel{B}\lambda}{\cancel{B} \left(-4n^{2} + \omega_{0}^{2}\right)} \Longrightarrow \textcolor{rojo}{\varphi_n = \arctan\left(\frac{2n\lambda}{-4n^{2} + \omega_{0}^{2}}\right)}
\end{gather*}
\begin{center}
    Usando $(2)^{2}$ + $(1)^{2}$:
\end{center}
\begin{align*}
    \left(\frac{4}{mn\pi}\right)^{2} \left(\cancel{\sin^{2}(\varphi) + \cos^{2}(\varphi)}\right) &= \left(2nB\lambda\right)^{2} + \left(B\left(-4n^{2} + \omega_{0}^{2}\right)\right)^{2}\\
    \frac{16}{m^{2}n^{2}\pi^{2}} &= 4n^{2}\lambda^{2} B^{2} + B^{2}\left(-4n^{2} + \omega_{0}^{2}\right)^{2}\\
    \frac{16}{m^{2}n^{2}\pi^{2}} &= B^{2} \left[4n^{2}\lambda^{2} + \left(-4n^{2} + \omega_{0}^{2}\right)^{2}\right]\\
    B^{2} &= \left(\frac{16}{m^{2}n^{2}\pi^{2}}\right) \left(\frac{1}{4n^{2}\lambda^{2} + \left(-4n^{2} + \omega_{0}^{2}\right)^{2}}\right)\\
    B &= \frac{4}{mn\pi} \left(\frac{1}{\sqrt{4n^{2}\lambda^{2} + \left(-4n^{2} + \omega_{0}^{2}\right)^{2}}}\right)
\end{align*}
\begin{center}
    Para este caso, $B$ puede ser reescrito como $B_n$ siendo netamente función de $n$, por tanto:
\end{center}
\begin{gather*}
    B = m\frac{\pi}{4} B_n\\
    \textup{donde,}\\
    \textcolor{rojo}{B_n = n\left(4n^{2}\lambda^{2} + \left(-4n^{2} + \omega_{0}^{2}\right)^{2}\right)^{-\sfrac{1}{2}}}
\end{gather*}

\centerline{\rule{0.95\textwidth}{0.4pt}} % Línea negra

\begin{multicols}{2}

Tras realizado el tratamiento mostrado con anterioridad, tenemos los valores para $B$ y para $\varphi$, conocemos su naturaleza como constantes ya que no dependen en ningún momento de un parámetro como $t$, aunque dependan de $n$ siendo funciones de ella, recordando la naturaleza de $n$ podemos seguir diciendo que $B$ y $\varphi$ son constantes.

Usando los resultados en rojo para $B$ y $\varphi$, transformadas en función de $n$ y teniendo que $A = 0$, la solución particular $Z_p$ es:
\begin{equation*}
    Z_p = \sum_{\mathrm{impar}}^{\infty} m\frac{\pi}{4} B_n e^{i\left(2nt - \frac{\pi}{2} - \varphi_n\right)}
\end{equation*}
Pero, en realidad lo que necesitamos es la parte real de esa expresión que por tratamientos previamente hechos ya sabemos que la obtenemos desarrollando la exponencial compleja, de esta manera volviendo a la forma original para $F(t)$, teniendo así la solución particular $x_p$, resulta:
\begin{equation}
    \textcolor{rojo}{x_p = m\frac{\pi}{4} \sum_{\mathrm{impar}}^{\infty} B_n \sin\left(2nt - \varphi_n\right)} \label{pro5:res4}
\end{equation}

$b)$ Para el caso 2, con:
\begin{equation*}
    F(t) =
    \begin{cases}
    1 + \frac{4t}{\pi},&\text{si } \sfrac{-\pi}{2} \leq t < 0\\
    1 - \frac{4t}{\pi},&\text{si } \hspace{18pt} 0 \leq t < \sfrac{\pi}{2}
    \end{cases}
\end{equation*}
\begin{center}$F(t)=F(t+\pi)$\end{center}

Gracias al caso anterior, podemos sustentarnos de esa misma explicación para este caso volviendo así mucho más corto el desarrollo, sin embargo, no tenemos la serie de Fourier para este ejercicio ya hecha por lo que es necesario calcularla, para ello nos vamos a ayudar del software \textit{DESMOS} de tal manera que conozcamos la naturaleza de la función, es decir, si es par o impar y así saber que coeficientes calcular.

\begin{figure}[H]
\begin{center}
\includegraphics[scale=0.28]{fig2.png}
\caption{Gráfica para $F(t)$ junto a su serie de Fourier. Puede verla aquí: \href{https://www.desmos.com/calculator/y9swhr0e0u}{\textit{Gráfica $F(t)$}}}
\label{fig2}
\end{center}
\end{figure}

En la Figura \ref{fig2} se puede ver la función $F(t)$ para este caso, junto a su serie de Fourier calculada, de ahí nos podemos dar cuenta que la serie pareciera ser primeramente par, sin embargo, el valor de $a_0$ es 0, corroborando los resultados su coeficiente $b_n$ también es igual a 0, teniendo así que a la hora de hallar el valor de $a_0$ la integral se anula. A continuación el tratamiento usando (\ref{pro2:eq1}):
\begin{align*}
    a_0 &= \frac{4}{\pi} \int\limits_{0}^{\sfrac{\pi}{2}} \left(1 - \frac{4t}{\pi}\right) \,dt\\
    &= \frac{4}{\pi} \left[\int\limits_{0}^{\sfrac{\pi}{2}} \,dt - \frac{4}{\pi} \int\limits_{0}^{\sfrac{\pi}{2}} t \,dt\right]\\
    &= \frac{4}{\pi} \left[\left(\left. t \right|_{0}^{\sfrac{\pi}{2}}\right) - \frac{4}{\pi}\left(\left.\frac{t^{2}}{2}\right|_{0}^{\sfrac{\pi}{2}}\right)\right]\\
    &= \frac{4}{\pi} \left[\frac{\pi}{2} - \frac{\cancel{4}}{\cancel{\pi}}\left(\frac{\pi^{\cancel{2}}}{\cancel{8}}\right)\right]\\
    &= \frac{4}{\pi} \cancel{\left[\frac{\pi}{2} - \frac{\pi}{2}\right]}\\
    &= 0
\end{align*}
De esta manera, tenemos que $a_0 = b_0 = 0$, quedándonos solamente por hallar el valor de $a_n$. Usando (\ref{pro2:eq2}), resulta:
\begin{align*}
    a_n &= \frac{4}{\pi} \int\limits_{0}^{\sfrac{\pi}{2}} \left(1 - \frac{4t}{\pi}\right) \cos(2nt) \,dt\\
    &= \frac{4}{\pi} \left[\int\limits_{0}^{\sfrac{\pi}{2}} \cos(2nt) \,dt - \frac{4}{\pi} \int\limits_{0}^{\sfrac{\pi}{2}} t \cos(2nt) \,dt\right]
\end{align*}
Sustituyendo al usar integración por partes:
\begin{align*}
    u&=t          &  dv&=\cos\left(2nt\right) dt\\
    du&=dt         &  v&=\frac{1}{2n} \sin\left(2nt\right)
\end{align*}
Desarrollando esta integral por separado:
\begin{gather*}
    \int\limits_{0}^{\sfrac{\pi}{2}} t \cos(2nt) \,dt\\
    \left[\frac{t}{2n} \sin\left(2nt\right) - \frac{1}{2n} \int \sin\left(2nt\right) dt\right]_{0}^{\sfrac{\pi}{2}}\\
    \left[\frac{t}{2n} \sin\left(2nt\right) - \frac{1}{2n} \left(-\frac{1}{2n}\cos(2nt)\right)\right]_{0}^{\sfrac{\pi}{2}}\\
    \left[\frac{t}{2n} \sin\left(2nt\right) + \frac{1}{4n^{2}} \cos(2nt)\right]_{0}^{\sfrac{\pi}{2}}
\end{gather*}
Evaluando en los límites de integración y operando resulta:
\begin{equation*}
    \int\limits_{0}^{\sfrac{\pi}{2}} t \cos(2nt) \,dt = \frac{(-1)^{n}-1}{4n^{2}}
\end{equation*}
Reemplazando lo obtenido en el desarrollo del coeficiente:
\begin{align*}
    a_n &= \frac{4}{\pi} \left[\cancel{\left(\left.\frac{1}{2n}\sin(2nt)\right|_{0}^{\sfrac{\pi}{2}}\right)} - \frac{\cancel{4}}{\pi}\left(\frac{(-1)^{n}-1}{\cancel{4}n^{2}}\right)\right]\\
    &= \frac{4}{\pi} \left[-\frac{(-1)^{n}-1}{n^{2}\pi}\right]\\
    &= \frac{4\left(1-(-1)^{n}\right)}{n^{2}\pi^{2}}
\end{align*}
Teniendo el valor de $a_n$, la serie de Fourier para $F(t)$ es:
\begin{equation}
    F(t) = \sum_{n=1}^{\infty} \frac{4\left(1-(-1)^{n}\right)}{n^{2}\pi^{2}} \cos(2nt) \label{pro5:res5}
\end{equation}
Para este caso, no es necesario colocar la función de la forma usada en el caso anterior, ya que ese tratamiento se realiza por la facilidad de transformación del coseno como exponencial compleja. Por tanto podemos expresar directamente la parte compleja para $F(t)$ como:
\begin{equation}
    Z(t) = \sum_{n=1}^{\infty} \frac{4\left(1-(-1)^{n}\right)}{n^{2}\pi^{2}} e^{i2nt} \label{pro5:res6}
\end{equation}
Suponiendo el mismo anzats que del caso anterior y haciendo las derivadas, resulta:
\begin{align*}
    Z_p &= A + \sum_{n=1}^{\infty} Be^{i\left(2nt - \varphi\right)}\\
    \dot{Z_p} &= \sum_{n=1}^{\infty} i2nB e^{i\left(2nt - \varphi\right)}\\
    \ddot{Z_p} &= \sum_{n=1}^{\infty} -4n^{2}B^{2} e^{i\left(2nt - \varphi\right)}
\end{align*}
Reemplazando en la ecuación diferencial, desarrollando y usando (\ref{pro5:res6}):

\end{multicols}

\centerline{\rule{0.95\textwidth}{0.4pt}} % Línea negra

\begin{gather*}
    \sum_{n=1}^{\infty} -4n^{2}B^{2} e^{i\left(2nt - \varphi\right)} + \lambda\left(\sum_{n=1}^{\infty} i2nB e^{i\left(2nt - \varphi\right)}\right) + \omega_{0}^{2}\left(A + \sum_{n=1}^{\infty} B e^{i\left(2nt - \varphi\right)}\right) = \frac{1}{m} Z(t)\\
    \left[\sum_{n=1}^{\infty} e^{i\left(2nt\right)}e^{-\varphi} \left(-4n^{2}B^{2} + \lambda i2nB + \omega_{0}^{2} B\right) \right] + \omega_{0}^{2} A = \frac{1}{m} \left(\sum_{n=1}^{\infty} \frac{4\left(1-(-1)^{n}\right)}{n^{2}\pi^{2}} e^{i2nt}\right)
\end{gather*}
\begin{center}
    Si igualamos las componentes similares, obtenemos qué:
\end{center}
\begin{gather*}
    \omega_{0}^{2} A = 0 \Longrightarrow A = 0\\
    e^{-\varphi} \left(-4n^{2}B^{2} + \lambda i2nB + \omega_{0}^{2} B\right) = \frac{4\left(1-(-1)^{n}\right)}{mn^{2}\pi^{2}}
\end{gather*}
\begin{center}
    Desarrollando lo obtenido para $B$ y para $\varphi$, resulta:
\end{center}
\begin{gather*}
    \frac{4\left(1-(-1)^{n}\right)}{mn^{2}\pi^{2}} = e^{-\varphi} \left(-4n^{2}B^{2} + \lambda i2nB + \omega_{0}^{2} B\right)\\
    \frac{4\left(1-(-1)^{n}\right)}{mn^{2}\pi^{2}} e^{\varphi} = B\left(-4n^{2} + \omega_{0}^{2}\right) + i\left(2nB\lambda\right)\\
    \frac{4\left(1-(-1)^{n}\right)}{mn^{2}\pi^{2}} \left(\cos(\varphi) + i\sin(\varphi)\right) = B\left(-4n^{2} + \omega_{0}^{2}\right) + i\left(2nB\lambda\right)
\end{gather*}
\begin{center}
    Si se iguala la parte real con la compleja en ambas partes, se genera el siguiente sistema de ecuaciones:
\end{center}
\begin{gather*}
    \begin{cases}
    \frac{4\left(1-(-1)^{n}\right)}{mn^{2}\pi^{2}} \cos(\varphi) = B\left(-4n^{2} + \omega_{0}^{2}\right) & (1)\\
    \frac{4\left(1-(-1)^{n}\right)}{mn^{2}\pi^{2}} \sin(\varphi) = 2nB\lambda & (2)
    \end{cases}
\end{gather*}
\begin{center}
    Usando (2) en (1):
\end{center}
\begin{gather*}
    \frac{\cancel{\frac{4\left(1-(-1)^{n}\right)}{mn^{2}\pi^{2}}}}{\cancel{\frac{4\left(1-(-1)^{n}\right)}{mn^{2}\pi^{2}}}} \tan(\varphi) = \frac{2n\cancel{B}\lambda}{\cancel{B} \left(-4n^{2} + \omega_{0}^{2}\right)} \Longrightarrow \textcolor{rojo}{\varphi_n = \arctan\left(\frac{2n\lambda}{-4n^{2} + \omega_{0}^{2}}\right)}
\end{gather*}
\begin{center}
    Usando $(2)^{2}$ + $(1)^{2}$:
\end{center}
\begin{align*}
    \left(\frac{4\left(1-(-1)^{n}\right)}{mn^{2}\pi^{2}}\right)^{2} \left(\cancel{\sin^{2}(\varphi) + \cos^{2}(\varphi)}\right) &= \left(2nB\lambda\right)^{2} + \left(B\left(-4n^{2} + \omega_{0}^{2}\right)\right)^{2}\\
    \frac{16\left(1-(-1)^{n}\right)^{2}}{m^{2}n^{4}\pi^{4}} &= 4n^{2}\lambda^{2} B^{2} + B^{2}\left(-4n^{2} + \omega_{0}^{2}\right)^{2}\\
    \frac{16\left(1-(-1)^{n}\right)^{2}}{m^{2}n^{4}\pi^{4}} &= B^{2} \left[4n^{2}\lambda^{2} + \left(-4n^{2} + \omega_{0}^{2}\right)^{2}\right]\\
    B^{2} &= \left(\frac{16\left(1-(-1)^{n}\right)^{2}}{m^{2}n^{4}\pi^{4}}\right) \left(\frac{1}{4n^{2}\lambda^{2} + \left(-4n^{2} + \omega_{0}^{2}\right)^{2}}\right)\\
    B &= \frac{4\left(1-(-1)^{n}\right)}{mn^{2}\pi^{2}} \left(\frac{1}{\sqrt{4n^{2}\lambda^{2} + \left(-4n^{2} + \omega_{0}^{2}\right)^{2}}}\right)
\end{align*}
\begin{center}
    Para este caso, $B$ puede ser reescrito como $B_n$ siendo netamente función de $n$, por tanto:
\end{center}
\begin{gather*}
    B = m\frac{\pi^{2}}{4} B_n
\end{gather*}\
\begin{center}
    donde,
\end{center}
\begin{gather*}
    \textcolor{rojo}{B_n = \frac{1-(-1)^{n}}{n^{2} \sqrt{4n^{2}\lambda^{2} + \left(-4n^{2} + \omega_{0}^{2}\right)^{2}}}}
\end{gather*}

\centerline{\rule{0.95\textwidth}{0.4pt}} % Línea negra

\begin{multicols}{2}

Tras realizado el tratamiento mostrado, tenemos los valores para $B$ y $\varphi$, nuevamente constantes y sus transformadas en función netamente de $n$. De esta manera, reemplazando para $Z_p$:
\begin{equation*}
    Z_p = \sum_{n=1}^{\infty} m\frac{\pi^{2}}{4} B_n e^{i\left(2nt - \varphi_n\right)}
\end{equation*}
Pero, como necesitamos es su parte real, se desarrolla la exponencial compleja lo que resulta:
\begin{equation}
    \textcolor{rojo}{x_p = m\frac{\pi^{2}}{4} \sum_{n=1}^{\infty} B_n \cos(2nt - \varphi_n)} \label{pro5:res7}
\end{equation}

Finalizando con el ejercicio, a continuación se expresan ambas soluciones generales para las ecuaciones diferenciales no homogéneas para cada caso.

\end{multicols}

\centerline{\rule{0.95\textwidth}{0.4pt}} % Línea negra

\subsubsection*{Caso 1:}

\begin{equation*}
    \frac{d^{2}x}{dt^{2}} + \lambda \frac{dx}{dt} + \omega_{0}^{2}x = \frac{1}{m}F(t) \textup{ con }
    F(t) = F(t+\pi) =
    \begin{cases}
        -1,&\text{si } \sfrac{-\pi}{2} \leq t < 0\\
        1,&\text{si } \hspace{18pt} 0 \leq t < \sfrac{\pi}{2}
    \end{cases}
\end{equation*}

La solución de la ecuación diferencial es:

\begin{equation}
    \textcolor{rojo}{x(t) = A_0e^{-\frac{\lambda}{2}} \cos\left(\omega_a t + \alpha_a\right) + m\frac{\pi}{4} \sum_{\mathrm{impar}}^{\infty} B_n \sin\left(2nt - \varphi_n\right)} \label{pro5:resf1}
\end{equation}

Donde:

\begin{equation*}
    B_n = n\left(4n^{2}\lambda^{2} + \left(-4n^{2} + \omega_{0}^{2}\right)^{2}\right)^{-\sfrac{1}{2}} \textup{, y }\varphi_n = \arctan\left(\frac{2n\lambda}{-4n^{2} + \omega_{0}^{2}}\right)
\end{equation*}

\subsubsection*{Caso 2:}

\begin{equation*}
    \frac{d^{2}x}{dt^{2}} + \lambda \frac{dx}{dt} + \omega_{0}^{2}x = \frac{1}{m}F(t) \textup{ con }
    F(t) = F(t+\pi) =
    \begin{cases}
        1 + \frac{4t}{\pi},&\text{si } \sfrac{-\pi}{2} \leq t < 0\\
        1 - \frac{4t}{\pi},&\text{si } \hspace{18pt} 0 \leq t < \sfrac{\pi}{2}
    \end{cases}
\end{equation*}

La solución de la ecuación diferencial es:

\begin{equation}
    \textcolor{rojo}{x(t) = A_0e^{-\frac{\lambda}{2}} \cos\left(\omega_a t + \alpha_a\right) + m\frac{\pi^{2}}{4} \sum_{n=1}^{\infty} B_n \cos(2nt - \varphi_n)} \label{pro5:resf2}
\end{equation}

Donde:

\begin{equation*}
    B_n = \frac{1-(-1)^{n}}{n^{2} \sqrt{4n^{2}\lambda^{2} + \left(-4n^{2} + \omega_{0}^{2}\right)^{2}}} \textup{, y }\varphi_n = \arctan\left(\frac{2n\lambda}{-4n^{2} + \omega_{0}^{2}}\right)
\end{equation*}

\centerline{\rule{0.95\textwidth}{0.4pt}} % Línea negra

\begin{multicols}{2}

\begin{enumerate}[leftmargin=15pt]
    \item[$2$.] Averiguar como se calcula la transformada de Laplace de la función periódica $F(t)$ de periodo $T_{F}$, tal que $F(t) =F(t+T_{F})$. Luego calcule la Transformada de Laplace de la función periódica  del caso 1 del problema anterior, $\mathscr{L}\{F(t)\}$:
    \[ F(t) = 
      \left \{
        \begin{aligned}
          -1 &,\ \text{si} \ -\pi/2 \le t < 0\\
          1 &,\ \text{si} \ 0 \le t < \pi/2 \\
        \end{aligned}
      \right .
    \] 
    \[y\]
    \[
    F(t) = F(t+\pi)\]
    
    ¿Cómo se resolvería la ecuación diferencial  \[\frac{d^2x}{dt^2}+\lambda\frac{dx}{dt}+\omega_{0}^2x= \frac{1}{m}F(t) \] usando la técnica de la transformada de Laplace?. ¡Realice una propuesta!
\end{enumerate}
    
\end{multicols}

\end{document}